\chapter{Representation Invariance}
\label{ch:invariance}

\section{Statement}

The question addressed in this chapter is when two representations,
possibly realized in entirely different substrates, ought to be
considered the same object. The answer given here is that sameness of
substrate is irrelevant; what matters is sameness of admissibility
structure.

\begin{definition}[Admissibility-isomorphism]
Let $R$ be an admissibility structure on a set $X$ and $R'$ an
admissibility structure on a set $X'$. A map $\varphi : X \to X'$ is an
\emph{admissibility-isomorphism} if, for all $x, y \in X$,
\[
x \adm{R} y \quad \iff \quad \varphi(x) \adm{R'} \varphi(y).
\]
\end{definition}

\begin{theorem}[Representation Invariance]
\label{thm:invariance}
If there exists an admissibility-isomorphism $\varphi : X \to X'$
between $(X, R)$ and $(X', R')$, then $(X, R)$ and $(X', R')$ are
semantically equivalent, regardless of whatever physical or
computational substrate realizes $X$ and $X'$.
\end{theorem}

\section{Discussion}

The theorem as stated is close to definitional, and this is by design
rather than by accident. Semantic equivalence, within this research
program, has never been defined as identity of physical encoding; it has
been defined, throughout the treatment of admissible representation and
context preservation developed elsewhere in this program, as identity of
what a representation permits and forbids as continuation. Theorem
\ref{thm:invariance} simply makes that definition explicit as an
invariance statement rather than leaving it distributed across the
apparatus that depends on it.

What is not definitional, and what does real work when the theorem is
applied, is establishing that a particular $\varphi$ is in fact an
admissibility-isomorphism between two representations of interest. This
is generally the harder direction: showing that a biphasic transformation
--- structured meaning realized as continuous speech, or as sequential
text, or as embodied gesture --- preserves the reachable-set structure
exactly, rather than merely preserving some coarser statistic that
happens to look similar under casual inspection. The content of a
representation-invariance claim, in any particular application, is
carried entirely by the verification that $\varphi$ satisfies the
biconditional in both directions: that admissible continuations in one
representation correspond exactly to admissible continuations in the
other, with neither representation admitting a continuation the other
forbids.

\begin{remark}
It follows immediately from Theorem \ref{thm:invariance} that
representation invariance is an equivalence relation on admissibility
structures: the identity map witnesses reflexivity, the inverse of an
admissibility-isomorphism is again an admissibility-isomorphism, giving
symmetry, and composition of admissibility-isomorphisms is again an
admissibility-isomorphism, giving transitivity. Consequently, the class
of representations equivalent to a given $(X, R)$ partitions into a
single equivalence class under $\varphi$, and asking whether two
representations are "the same" reduces to asking whether they lie in the
same class.
\end{remark}

\begin{corollary}
\label{cor:invariance-substrate}
No property of $X$ or $X'$ that is not preserved by every
admissibility-isomorphism can be a semantic property of the
representation. In particular, substrate-specific properties --- the
particular physical channel through which a representation is realized,
its encoding cost, or its production latency --- are not semantic
properties under this framework unless they happen to be expressible as
constraints on $R$ itself.
\end{corollary}

Corollary \ref{cor:invariance-substrate} is worth stating explicitly
because it draws a sharp boundary that is easy to blur in practice.
Structural semantics, as developed elsewhere in this program, is
precisely the study of what survives this partition; anything that does
not survive it belongs to a different, and legitimate, subject ---
the study of channel, cost, and implementation --- but not to the study
of what a representation \emph{means}.
